% =====================================================================
% DRAFT FIX for the factorized-state definition (Section 4).
% NOT IN THE BUILD. Review, then merge into elsarticle-template-num.tex.
%
% THE DEFECT
% The manuscript declares  z* = [z^v, z^r] in R^6  with
%     z^v = (x_rel, y_rel, psi_rel)      and      z^r = (e_CTE, e_psi, kappa).
% Four things are wrong, all checkable against src/models/frenet_dynamics.py:
%
%   1. v and omega are propagated by eqs. (16)-(19) and by the implementation
%      (frenet_dynamics.py:103-104), but appear in neither component.
%   2. The arc length s is propagated (line 100) and is what kappa is indexed by,
%      but appears nowhere in the declared state.
%   3. kappa is declared AS a state component but is never propagated: it is read,
%      either from the map (kappa_at(s), line 82) or from the perceived preview.
%      This is the paper's own central claim -- the road is re-indexed, not
%      predicted -- so having kappa in the propagated state contradicts it.
%   4. (x_rel, y_rel, psi_rel) is the conference/v5 parameterisation. The Frenet
%      model replaces it with (s, d, psi_e); it does not carry both.
%
% Root cause: Section 4 is v5-era text that was never reconciled with the Frenet
% formulation introduced in Section 5.
%
% THE ONE JUDGEMENT CALL FOR THE AUTHOR
% The rewrite below keeps the z^v / z^r split, because the surrounding narrative
% and Fig. 2 are built on it, and maps it onto the implemented state. The
% alternative is to drop the split and present a single five-dimensional Frenet
% state with the road as a separate perceived function. That is arguably cleaner
% but touches more of Sections 4-5 and the architecture figure.
% =====================================================================

\paragraph{Vehicle state}
In the road-aligned frame the vehicle's configuration is given by its arc length
along the reference centerline, its speed and its yaw rate:
\begin{equation}
    \zv_t =
    \begin{bmatrix}
        s_t \\ v_t \\ \omega_t
    \end{bmatrix},
    \label{eq:vehicle_rel}
\end{equation}
where $s_t$ is the distance travelled along the centerline, $v_t$ the
longitudinal speed and $\omega_t$ the yaw rate. Unlike the global-pose
parameterization used in the fully observable setting, $s_t$ locates the vehicle
\emph{along the road} rather than in the world, which is what allows the road
description below to be indexed by it.

\paragraph{Road-context state}
The road-context component $\zr_t$ captures the vehicle's relationship to the
road geometry:
\begin{equation}
    \zr_t =
    \begin{bmatrix}
        e_{\mathrm{CTE},t} \\ e_{\psi,t}
    \end{bmatrix},
    \label{eq:road_ctx}
\end{equation}
where $e_{\mathrm{CTE}}$ is the cross-track error (signed lateral distance from
the lane center) and $e_{\psi}$ is the heading error (difference between the
vehicle's heading and the road tangent at its current position).

\paragraph{Factorized state}
The complete \emph{factorized interpretable state} is
\begin{equation}
    \zstar_t = \left[ \zv_t,\; \zr_t \right]
             = \left( s_t,\, v_t,\, \omega_t,\, e_{\mathrm{CTE},t},\, e_{\psi,t} \right)
             \in \mathbb{R}^5 .
    \label{eq:factorized_state}
\end{equation}

\paragraph{The road is not part of the state}
The local curvature $\kappa$ is deliberately \emph{absent} from
\eqref{eq:factorized_state}. It is not propagated. Instead the model carries a
curvature profile $\hat\kappa(\cdot)$ over arc length, perceived once from the
image window at the start of the rollout, and reads the curvature at the
vehicle from it:
\begin{equation}
    \kappa_t = \hat\kappa\!\left(s_t - s_{t_0}\right).
    \label{eq:kappa_readout}
\end{equation}
This is the distinction the rest of the article rests on. A model that carries
curvature as a state variable must predict how it evolves, and that prediction
compounds over a long rollout. Re-indexing a profile that was observed once
cannot compound, because nothing about the road is being extrapolated --- only
the vehicle's position along it advances.

Consequently the Markov property holds for the pair
$\left(\zstar_t,\, \hat\kappa(\cdot)\right)$, with $\hat\kappa$ held fixed
across the rollout: given the current factorized state, the perceived profile
and the action, the next factorized state is determined up to the unknown
parameters $\theta^*$. It does \emph{not} hold for $\zstar_t$ alone, since
\eqref{eq:kappa_readout} requires the profile. We state the condition in this
form because it is the one the implementation satisfies, and because the fixed
profile is exactly what bounds the long-horizon error.

% ---------------------------------------------------------------------
% KNOCK-ON EDITS REQUIRED ELSEWHERE (do not merge the above without these)
%
%   * Any prose describing z* as six-dimensional or listing kappa among the state
%     components.
%   * Figure 2's caption, if it enumerates the state components.
%   * The claim that the two branches' "only coupling" is \hat{z*}_{t+1}: the
%     actuation map inside phi_theta consumes z^r_t (frenet_dynamics.py:92-95
%     feeds d, psi, v, omega and kappa into dv_net/dom_net). Either restate the
%     coupling or drop the "only".
%   * Figure 7(b) perturbs v and omega as state variables. Under this definition
%     that is now consistent -- previously it was not.
% =====================================================================
